\section{Conclusion}


本稿は、意識を内部で生成される対象としてではなく、異なる時間スケールの処理のあいだに開かれるログ参照状態として再定式化した。

中心となる主張は単純である。

本稿では、意識を生成された対象としてモデル化しない。

意識は、高速な予測過程と低速なログ沈殿過程が、完全同期にも完全分離にも至らない範囲で結合したときに開かれる。

この開口には、速度差 $\Delta v$、結合粘性 $\mu$、内部遅延 $\tau$、意味勾配 $\nabla \Phi$、有効トルク様量 $\mathcal{T}_{\Phi}$ が関与する。

しかし、これらのいずれか一つが意識そのものなのではない。

意識は、それらが一定の関係に入ったときに成立する時間的非閉包レジームである。

\subsection{What Has Been Proposed}

本稿では、まず SSO を意識が成立するための最小時間構造として導入した。

SSO は、異なる速度で走る処理層が、中央制御なしに調整される構造である。

高速層は未来を先取りし、低速層は過去を沈殿させる。

そのあいだに有限の速度差が残ることで、現在処理と過去ログのあいだに参照関係が開かれる。

次に、トルクコンバータ・アテンションを導入した。

これは、注意を選択や重みづけとしてではなく、速度差を破綻なく変換する時間的機構として扱うものである。

この立場では、自我は主体ではなく、高速層と低速層の速度差が変換される位置として理解される。

さらに、IFGT と VED を用いて、ログ参照を意味勾配と非閉包ダイナミクスの中に再配置した。

IFGT は、情報密度 $I$、情報流 $J_I$、局所的なログ参照流 $J_{\log}$、情報ポテンシャル $\Phi$ の射影的使用によって、どのログが参照されやすいかを記述する。

VED は、完全閉包が通常状態として実現されず、差と流れを保ち続ける構造を与える。

この二つを接続することで、意識は、意味勾配に方向づけられたログ参照が、時間的非閉包レジームの中で開かれる状態として記述された。

\subsection{What Has Been Repositioned}

本稿では、いくつかの重要な概念を再配置した。

第一に、意識は対象ではなく状態領域として再配置された。

意識は、どこかに保存される物体でも、脳内に生成される出力でもない。

それは、処理とログが参照可能な関係に入ったときに開かれるレジームである。

第二に、クオリアは意識の本体ではなく、沈殿境界のライブな前景として再配置された。

クオリアは否定されない。

しかし、それは保存される対象ではない。

保存されるのは、クオリアそのものではなく、クオリアが通過した後のログ構造である。

第三に、測定は意識をそのまま捕まえる操作ではなく、ライブな境界を沈殿残余へ変換する操作として再配置された。

測定は無意味ではない。

しかし、測定が捉えるのは、意識そのものではなく、意識が開かれた後に残る痕跡である。

第四に、非局在性は神秘化ではなく、意識を時間的関係として扱うことの帰結として再配置された。

意識には神経基盤がある。

しかし、意識状態そのものは単一部位に局在しない。

それは、複数の局所基盤が、特定の時間的・意味的・非閉包的関係に入ったときに開かれる。

第五に、病理や変性状態は、疾患分類ではなく、consciousness window からの逸脱方向として再配置された。

この整理により、意識状態は有無の二値ではなく、崩れ方をもつ動的レジームとして扱われる。

\subsection{The Central Formula}

本稿の中心構造は、次の式に要約できる。

\begin{equation}
\text{Consciousness}
\sim
\text{Log Referencing}
\left(
\Delta v,
\mu,
\tau,
\mathcal{T}_{\Phi},
\nabla \Phi
\right)
\end{equation}

この式は、意識を計算する式ではない。

それは、意識について語るときに必要な構造変数を示す座標である。

より具体的には、意識状態は次のような窓の中で安定する。

\begin{equation}
0 < |\Delta v| < \Delta v_{\max}
\end{equation}

\begin{equation}
\mu_{\min} < \mu < \mu_{\max}
\end{equation}

\begin{equation}
\tau_{\min} < \tau < \tau_{\max}
\end{equation}

\begin{equation}
\mathcal{T}_{\min}
<
\mathcal{T}_{\Phi}
<
\mathcal{T}_{\max}
\end{equation}

この consciousness window の内側では、処理は完全同期せず、完全分離もせず、ログ参照が安定して開かれる。

窓の外側では、意識状態はそれぞれ異なる方向へ変形する。

したがって、意識は点ではなく窓である。

対象ではなくレジームである。

生成物ではなく開口である。

\subsection{Why the Model Remains Open}

本稿は、意識の完全な統一方程式を与えない。

これは欠落ではなく、意図的な制約である。

意識を閉じた対象として定義しようとすると、意識を成立させている非閉包性そのものが失われる。

クオリアを保存物として扱えば、境界現象としての性質が失われる。

神経相関を意識そのものと同一視すれば、時間的レジームとしての構造が見えなくなる。

病理を診断名に直結させれば、consciousness window からの逸脱方向という柔軟な構造が失われる。

したがって、本稿はあえて閉じ切らない。

閉じ切らないことは、説明の放棄ではない。

本稿の主題である意識が、そもそも閉じ切らない時間構造として成立しているからである。

理論は対象に似る。

意識が非閉包的であるなら、意識理論もまた、完全閉包を急いではならない。

\subsection{Toward Further Work}

今後の課題は明確である。

第一に、$\mathcal{T}_{\Phi}$ を、神経活動、身体状態、主観報告、言語ログに接続する具体的な測定可能量へ近づける必要がある。

第二に、IFGT の $I$、$J_I$、$\Phi$ および局所流 $J_{\log}$ を、脳内データや行動データへ写像する方法を精密化する必要がある。

第三に、consciousness window を、睡眠、夢、麻酔、没入、解離、熟練運動、言語報告などの状態に対して検討する必要がある。

第四に、クオリアの沈殿残余を扱う実験設計を考える必要がある。

第五に、自己と言語と社会的相互作用が、ログ参照構造をどのように拡張するかを扱う必要がある。

この最後の課題は、意識編から自己編・知性編への橋である。

意識は、ログ参照の開口である。

自己は、その開口が反復された後に低速層へ沈殿する地形である。

知性は、その地形を用いて、未来予測、言語、社会的ログ、外部環境との相互作用へ展開する運動である。

\subsection{Final Statement}

意識は、どこかで作られるものではない。

意識は、速く流れる処理と、ゆっくり沈むログのあいだに、一時的に開かれる通路である。

そこには速度差がある。

遅延がある。

粘性がある。

意味勾配がある。

そして、閉じ切らない流れがある。

この通路が開いているあいだ、処理はただ流れるだけではなく、過去ログを参照し、自己の地形を更新し、世界との関係を再構成する。

この状態を、私たちは意識と呼んでいる。

したがって、本稿の結論は、冒頭の命題へ戻る。

本稿では、意識を生成された対象としてモデル化しない。

意識は、時間的非閉包の中で、ログ参照が開かれるときに成立する。
